\section*{第 IV 章　両側単純環}\label{kapitel-iv.-zweiseitig-einfache-ringe}

\subsection*{§ 18. 補題}\label{ein-hilfssatz}

$\mathsf K$ を体、$\mathfrak G$ を $\mathsf K$ の自己同型（それ自身への同型写像）からなる群とする。$\mathfrak G$ に属する各自己同型によって不変な $\mathsf K$ の元全体は $\mathsf K$ の部分体 $\mathsf P$ をなし、これを $\mathfrak G$ の\emph{不変体}という。

$\tau\ne0$ を $\mathsf K$ の任意の元とすると、$\mathsf K$ の元 $\alpha$ に $\tau^{-1}\alpha\tau$ を対応させることによって $\mathsf K$ の自己同型が得られる。群論の対応する用語にならって、このような自己同型を\emph{内部自己同型}という。一つの体 $\mathsf K$ の内部自己同型すべてからなる群の不変体 $\mathsf P$ は、$\mathsf K$ の中心である。

いま、体 $\mathsf K$、その自己同型群 $\mathfrak G$、およびその不変体 $\mathsf P$ が与えられているとする。

\[
\mathfrak M=x_1\mathsf P+\cdots+x_n\mathsf P
\]

を階数 $n$ の $\mathsf P$-加群とする。$\mathfrak G$ の元 $\gamma$ の $\mathfrak M$ の元への作用を

\[
\gamma\cdot x_i=x_i,\qquad\text{したがって}\qquad
\gamma\sum \varrho_i x_i=\sum \varrho_i x_i
\]

によって定める。すると $\mathfrak G$ は拡大加群

\[
\mathfrak M_{\mathsf K}=x_1\mathsf K+\cdots+x_n\mathsf K
\]

の自己同型群となり、$\mathfrak M_{\mathsf K}$ における $\mathfrak G$ の不変領域は $\mathfrak M$ となる。この仮定のもとで、次の定理が成り立つ。

\begingroup\itshape
\[
\mathfrak C=z_1\mathsf K+\cdots+z_r\mathsf K
\]

が $\mathfrak M_{\mathsf K}$ の部分加群であり、各 $\gamma\in\mathfrak G$ によってそれ自身へ移される（各元が固定される必要はない）ならば、$\mathfrak C$ は $\mathfrak M$ のある部分加群 $\mathfrak c$ を $\mathsf K$ によって拡大した加群である：

\[
\mathfrak C=\mathfrak c_{\mathsf K}.
\]
\endgroup

\emph{証明。} 1. 予備的注意として、$\mathfrak c=\mathfrak C\cap\mathfrak M$ でなければならない。一般に
\[
\mathfrak c=\mathfrak c_{\mathsf K}\cap\mathfrak M
\]
だからである。

$\mathfrak c\subseteqq\mathfrak c_{\mathsf K}\cap\mathfrak M$ は明らかである。逆の包含を示すため、$\mathfrak c$ の基底 $y_1,\ldots,y_r$ を取り、
\[
\mathfrak c=y_1\mathsf P+\cdots+y_r\mathsf P
\]
とする。このとき
\[
\mathfrak c_{\mathsf K}=y_1\mathsf K+\cdots+y_r\mathsf K.
\]
$\mathfrak c_{\mathsf K}\cap\mathfrak M$ の各元 $c$ は、$\mathfrak c_{\mathsf K}$ の元としての基底表示
\[
c=\sum_{i=1}^{r}y_i\varkappa_i
\]
と、$\mathfrak M$ の元としての基底表示
\[
c=\sum_{i=1}^{r}y_i\varrho_i+\sum_{j=r+1}^{n}y_j\varrho_j
\]
をもつ。ここで $y_{r+1},\ldots,y_n$ は、$y_1,\ldots,y_r$ を $\mathfrak M$ の基底に補う $\mathfrak M$ の元である（M. Z. § 6）。両方の表示を $\mathfrak M_{\mathsf K}$ における基底表示とみなせるので、それらは一致し、$\varkappa_i=\varrho_i$（$i=1,\ldots,r$）、$\varrho_j=0$（$j=r+1,\ldots,n$）である。したがって
\[
c=\sum_{i=1}^{r}y_i\varrho_i,
\qquad\text{証明された。}
\]

\textbf{2.} 定理を証明するため、元 $z_i$ が $\mathfrak M$ に属するような $\mathfrak C$ の基底 $z_1,\ldots,z_r$ を構成する。このとき
\[
\mathfrak c=z_1\mathsf P+\cdots+z_r\mathsf P
\]
と置けば、実際に
\[
\mathfrak C=\mathfrak c_{\mathsf K}
\]
となる。

$\mathfrak C$ の任意の基底 $y_1,\ldots,y_r$ を取り、$\mathfrak M$ の一つの基底 $x_1,\ldots,x_n$ の適当な元 $x_{r+1},\ldots,x_n$ を付け加えて、$\mathfrak M_{\mathsf K}$ の基底にする。この基底は同時に $\mathfrak M_{\mathsf K}$ の基底でもある。この選択は M. Z. § 5 により可能である：
\[
\mathfrak M_{\mathsf K}=\mathfrak C+x_{r+1}\mathsf K+\cdots+x_n\mathsf K.
\]
ここから、$\mathfrak M$ の基底の最初の $r$ 元 $x_1,\ldots,x_r$ について
\[
x_i=z_i-\sum_{j=r+1}^{n}x_j\varkappa_j^{(i)},\qquad
z_i\in\mathfrak C,\quad \varkappa_j^{(i)}\in\mathsf K
\]
が得られる。$z_i$ は $x_{r+1},\ldots,x_n$ とともに明らかに $\mathfrak M_{\mathsf K}$ の基底をなすので線型独立であり、$z_1,\ldots,z_r$ は $\mathfrak C$ の基底である。この基底が求める形であること、すなわち $z_i\in\mathfrak M$ を示す。

$\gamma$ を $\mathfrak G$ の元とする。一方では
\[
\begin{array}{rcl}
(1)\qquad \gamma z_i
&=&\gamma x_i+\displaystyle\sum_{j=r+1}^{n}\gamma x_j\cdot \gamma\varkappa_j^{(i)}
=x_i+\displaystyle\sum_{j=r+1}^{n}x_j\gamma\varkappa_j^{(i)}.
\end{array}
\]
他方では、$\mathfrak C$ がすべての $\gamma\in\mathfrak G$ を許すから、
\[
\begin{array}{rcl}
(2)\qquad \gamma z_i&=&\displaystyle\sum_{\mu=1}^{r}z_\mu\sigma_\mu^{(i)},
\qquad \sigma_\mu^{(i)}\in\mathsf K.
\end{array}
\]
(2) はさらに
\[
\begin{array}{rcl}
(3)\qquad \gamma z_i
&=&\displaystyle\sum_{\mu=1}^{r}x_\mu\sigma_\mu^{(i)}
+\displaystyle\sum_{j=r+1}^{n}x_j\sum_{\mu=1}^{r}\varkappa_j^{(\mu)}\sigma_\mu^{(i)}
\end{array}
\]
と書ける。(1) と (3) の係数を比較すると
\[
\sigma_j^{(i)}=
\begin{cases}
1 & i=j,\\
0 & i\ne j
\end{cases}
\]
を得る。したがって
\[
\gamma\cdot z_i=z_i.
\]
ゆえに $z_i$ は $\mathfrak G$ の不変領域、すなわち $\mathfrak M$ に属する。証明された。

\subsection*{§ 19. 係数体の非可換拡大体における両側単純超複素系の表現}\label{darstellungen-zweiseitig-einfacher-hyperkomplexer-systeme-in-nichtkommutativen-erweiterungskoerpern-ihrer-koeffizientenkoerper}

\[
\mathfrak S=x_1\mathsf P+\cdots+x_n\mathsf P
\]

を両側単純、したがって完全可約な超複素系とし、$\mathsf K$ を中心 $\mathsf P$ の体とする。

\textbf{定理 1.} \emph{$\mathfrak S_{\mathsf K}$ は両側単純である（したがって完全可約でもある）。}

\emph{証明。} 1. $\mathfrak S_{\mathsf K}$ の各両側イデアル $\mathfrak a$ は、$\mathsf K$ の内部自己同型群 $\mathfrak G$ に関して不変である。実際、零でない $\mathsf K$ の元 $\varkappa$ に対して $\varkappa^{-1}\mathfrak a\varkappa=\bar{\mathfrak a}$ と置けば、$\mathfrak a$ のイデアル性により $\bar{\mathfrak a}\subseteqq\mathfrak a$ である。しかし $\bar{\mathfrak a}$ は $\mathsf K$-加群として $\mathfrak a$ と同じ $\mathsf K$-階数をもつので、$\bar{\mathfrak a}=\mathfrak a$ でなければならない。

\textbf{2.} 直前の補題により、$\mathfrak S_{\mathsf K}$ の両側イデアル $\mathfrak a$ は $\mathfrak S$ の部分加群 $\mathfrak A$ の拡大イデアル、$\mathfrak a=\mathfrak A_{\mathsf K}$ でなければならない。$\mathfrak A=\mathfrak S\cap\mathfrak a$ であるから、
\[
\mathfrak S\mathfrak A\subseteqq\mathfrak S\mathfrak S\cap\mathfrak S\mathfrak a\subseteqq\mathfrak S\cap\mathfrak a=\mathfrak A
\]
および
\[
\mathfrak A\mathfrak S\subseteqq\mathfrak S\mathfrak S\cap\mathfrak a\mathfrak S\subseteqq\mathfrak S\cap\mathfrak a=\mathfrak A
\]
となり、$\mathfrak A$ は $\mathfrak S$ の両側イデアルである。ところが $\mathfrak S$ の両側イデアルは $0$ と $\mathfrak S$ だけなので、$\mathfrak S_{\mathsf K}$ の両側イデアルも $0$ と $\mathfrak S_{\mathsf K}$ だけである。証明された。

次に $\mathfrak S$ の $\mathsf K$ における既約表現を求める。形式上の利点があるので、相反表現を構成する。直接表現と相反表現は一対一に対応するため、どちらを取っても結局は同じである。

$\mathfrak M$ を $\mathfrak S$ の $\mathsf K$ における相反表現加群、すなわち有限階数の $\mathfrak S$-左加群かつ $\mathsf K$-左加群
\[
\mathfrak M=\mathsf K\cdot y_1+\cdots+\mathsf K\cdot y_r
\]
で、交換則
\[
s(\varkappa m)=\varkappa (sm)
\]
を満たすものとする。

$\mathfrak M$ が $\mathfrak S_{\mathsf K}$-左加群ともみなせることを示す。まず $\mathfrak S_{\mathsf K}$ の元と $\mathfrak M$ の元との積を定めなければならない。$s\in\mathfrak S$、$x\in\mathsf K$、$m\in\mathfrak M$ に対して
\[
s\varkappa \cdot m=s\cdot(\varkappa \cdot m)
\]
と置き、一般の $\mathfrak S_{\mathsf K}$ の元 $\sum s_i\varkappa _i$ には
\[
\sum s_i\varkappa _i\cdot m=\sum s_i\cdot(\varkappa _i\cdot m)
\]
と定める。

この定義は一意であり、$\mathfrak S_{\mathsf K}$ の元を積 $s_i\varkappa _i$ の和として表す仕方に依存しない。すなわち
\[
\begin{array}{rcl}
(1)\qquad \displaystyle\sum_{i=1}^{p}s_i\varkappa _i&=&0
\end{array}
\]
から
\[
\sum_{i=1}^{p}s_i(\varkappa _i m)=0
\]
が従うことを示せばよい。$s_1,\ldots,s_q$ が線型独立で、$s_{q+1},\ldots,s_p$ が
\[
s_j=\sum_{i=1}^{q}s_i\varrho_i^{(j)},\qquad j=q+1,\ldots,p
\]
と表されるとする。表示は一意で、しかもすでに $\mathfrak S$ の中で可能だから、係数は $\mathsf P$ に属する。(1) から
\[
\varkappa _i+\sum_{j=q+1}^{p}\varrho_i^{(j)}\varkappa _j=0,\qquad i=1,\ldots,q
\]
を得る。したがって
\[
\begin{aligned}
\sum_{i=1}^{p}s_i(\varkappa _i m)
&=s_1(\varkappa _1m)+\cdots+s_q(\varkappa _qm)
  +\left(\sum_i s_i\varrho_i^{(q+1)}\right)(\varkappa _{q+1}m)\\
&\quad+\cdots+\left(\sum_i s_i\varrho_i^{(p)}\right)(\varkappa _pm)\\
&=s_1(\varkappa _1m)+\cdots+s_q(\varkappa _qm)\\
&\quad+s_1\varrho_1^{(q+1)}(\varkappa _{q+1}m)+\cdots+s_q\varrho_q^{(q+1)}(\varkappa _{q+1}m)
  \qquad\text{加群の性質により}\\
&\quad+\cdots\\
&\quad\cdots\\
&\quad+s_1\varrho_1^{(p)}(\varkappa _pm)+\cdots+s_q\varrho_q^{(p)}(\varkappa _pm)\\
&=s_1(\varkappa _1m)+\cdots+s_q(\varkappa _qm)\\
&\quad+s_1(\varrho_1^{(q+1)}\varkappa _{q+1}m)+\cdots+s_q(\varrho_q^{(q+1)}\varkappa _{q+1}m)\\
&\quad+\cdots\\
&\quad\cdots\\
&\quad+s_1(\varrho_1^{(p)}\varkappa _pm)+\cdots+s_q(\varrho_q^{(p)}\varkappa _pm)
  \qquad\begin{gathered}\text{結合則 }s'(sm)=s's\cdot m\text{ および}\\ \varkappa'(\varkappa m)=\varkappa'\varkappa\cdot m\text{ により}\end{gathered}\\
&=s_1\left(\varkappa _1m+\sum\varrho_1^{(q+1)}\varkappa _{q+1}\cdot m\right)+\cdots\\
&=s_1\left(\left(\varkappa _1+\sum\varrho_1^{(q+1)}\varkappa _{q+1}\right)m\right)+\cdots\\
&=0.
\end{aligned}
\]

$\mathfrak M$ の $\mathfrak S_{\mathsf K}$-加群としての分配則 $(S+S')m=Sm+S'm$、$S(m+m')=Sm+Sm'$ は明らかである。注意を要するのは結合則 $S(S'm)=SS'\cdot m$ である。分配則が成り立つので、
\[
s\varkappa \cdot(s'\varkappa ' m)=s\varkappa (s'\varkappa '\cdot m)
\]
の場合だけを考えればよい。一方、$s\varkappa \cdot m$ の定義により
\[
s\varkappa \cdot(s'\varkappa '\cdot m)=s\cdot(\varkappa \cdot(s'\cdot(\varkappa '\cdot m)))
\]
である。他方、
\[
s\varkappa  s'\varkappa '\cdot m=ss'\varkappa \varkappa '\cdot m=ss'\cdot(\varkappa \varkappa '\cdot m)
=s\cdot(s'\cdot(\varkappa \cdot(\varkappa '\cdot m)))
\]
は、二つの交換関係
\[
\varkappa \cdot s'=s'\cdot \varkappa
\]
および
\[
s'\cdot(\varkappa \cdot m)=\varkappa \cdot(s'\cdot m)
\]
により $s\cdot(x\cdot(s'\cdot(x'\cdot m)))$ に等しい。

したがって $\mathfrak M$ は実際に有限 $\mathfrak S_{\mathsf K}$-左加群である。逆に、各有限 $\mathfrak S_{\mathsf K}$-左加群は $\mathfrak S$ の $\mathsf K$ における相反表現加群である。実際、明らかに有限 $\mathsf K$-加群であり、
\[
s(\varkappa \cdot m)=s\varkappa \cdot m=\varkappa s\cdot m=\varkappa (s\cdot m)
\]
により交換条件も満たす（§ 24 のより簡単な証明も参照）。

M. Z. § 18 により、各有限 $\mathfrak S_{\mathsf K}$-加群は、$\mathfrak S_{\mathsf K}$ の単純左イデアルと $\mathfrak S_{\mathsf K}$-作用素同型な単純加群の和である。逆に $\mathfrak S_{\mathsf K}$ の各左イデアルは $\mathfrak S$ の $\mathsf K$ における表現加群であり、$\mathfrak S_{\mathsf K}$ の単純左イデアルはすべて作用素同型なので、次が成り立つ。

\textbf{定理 2.} \emph{$\mathfrak S$ の $\mathsf K$ における既約相反表現類はただ一つであり、$\mathfrak S_{\mathsf K}$ の一つの単純左イデアルによって生成される。}

この既約表現の次数を $r$ とすると、任意の表現の次数は $r$ の倍数である。もちろん既約表現を並べて結合すれば、$r$ の任意の倍数 $rs$ に対して $rs$ 次表現を構成できる。これは $s$ 個の単純加群の直和である表現加群を作ることに対応する。

$\mathfrak S$ 自身を表現加群として生成される $\mathfrak S$ の $\mathsf K$ における表現（主表現、M. Z. § 20）はすでに知られており、これは $\mathsf P$ においてすでに存在する。その次数は $n$ なので、$n$ は $r$ で割り切れる：
\[
n=rt.
\]
商 $t$ は容易に定まる。$r$ は一つの左イデアルの $\mathsf K$ 上の階数であり、$n$ は $\mathfrak S_{\mathsf K}$ 全体の階数だから、$t$ は $\mathfrak S_{\mathsf K}$ の左イデアルの個数である。したがって $t^2$ は、その単純イデアルの自己同型体上の $\mathfrak S_{\mathsf K}$ の階数である。

以下、$\mathfrak S$ が可換体である場合に特殊化する。

\textbf{定理 3.} \emph{$\mathfrak S$ が可換体ならば、$\mathfrak S$ は $\mathfrak S_{\mathsf K}$ の中心である。}

\emph{証明。} $\mathfrak S$ の各元は $\mathfrak S$ の各元および $\mathsf K$ の各元と、したがって $\mathfrak S_{\mathsf K}$ の各元と可換なので、$\mathfrak S$ は $\mathfrak S_{\mathsf K}$ の中心に属する。逆に、$\mathfrak S_{\mathsf K}$ の元がそのすべての元と可換であるならば、少なくとも $\mathsf K$ の各元と可換でなければならない。したがってその係数は $\mathsf K$ の中心 $\mathsf P$ に属し、その元は $\mathfrak S$ に属する。証明された。

可換環では相反同型と直接同型に差がないので、各相反表現は直接表現ともみなせる。したがって $\mathfrak S$ の $\mathsf K$ における既約表現は、直接表現加群によっても与えられなければならない。この直接表現加群は $\mathfrak S_{\mathsf K}$ の単純右イデアルである。実際、$\mathfrak r$ は階数 $r$ の有限 $\mathsf K$-右加群であり、$\mathfrak r\mathfrak S=\mathfrak S\mathfrak r$ により $\mathfrak S$-左加群でもあり（もちろん混合結合則を満たす）、階数 $r$ の直接表現加群である。既約表現類は一つだけなので、それは与えられた表現を生成しなければならない。

\subsection*{§ 20. 非可換体}\label{nichtkommutative-koerper}

ここで
\[
\mathfrak K=y_1\mathsf P+\cdots+y_m\mathsf P
\]
を、中心が $\mathsf P$ である非可換体とする。$\mathfrak K$ について、可換な有限代数体拡大の理論における諸定理と類似する一連の定理を展開する。

$\mathsf{Z}=\xi_1\mathsf P+\cdots+\xi_n\mathsf P$ を $\mathsf P$ 上の超複素系とすると、二つの拡大系 $\mathfrak K_{\mathsf{Z}}$ と $\mathsf{Z}_{\mathfrak K}$ は一致する。実際、
\[
\mathfrak K_{\mathsf{Z}}=y_1\xi_1\mathsf P+\cdots+y_m\xi_n\mathsf P
\]
および
\[
\mathsf{Z}_{\mathfrak K}=\xi_1y_1\mathsf P+\cdots+\xi_ny_m\mathsf P
\]
と書け、$y_i\xi_j=\xi_jy_i$ により両者は一致する。

以下では $\mathsf{Z}$ を常に可換体とするので、§ 19 の定理を適用できる。§ 19 の定理 1 により $\mathfrak K_{\mathsf{Z}}$ は常に両側単純である。

\textbf{定理 1.} \emph{$\Omega$ を $\mathsf P$ の代数閉拡大体とする。このとき $\mathfrak K_\Omega$ の中心は $\Omega$ に等しい。}

\emph{証明。} § 19 の定理 3 を用いる。$\Omega$ が $\mathfrak K_\Omega$ の中心に含まれることは明らかである。$\alpha$ を $\mathfrak K_\Omega$ の中心の任意の元とすると、その係数は $\mathsf P$ のある有限代数拡大体 $Z$（係数が生成する体）に属する。$\alpha$ はもちろん $\mathfrak K_{\mathsf{Z}}$ の中心にも属する。ところがその中心は $\mathsf{Z}$ なので、$\alpha\in Z$、したがって $\alpha\in\Omega$ である。証明された。

全く同様に次を証明する。

\textbf{定理 2.} \emph{$\mathfrak K_\Omega$ は両側単純である。}

\emph{証明。} $\mathfrak a=z_1\Omega+\cdots+z_r\Omega$ を両側イデアルとする。$z_1,\ldots,z_r$ の係数は $\mathsf P$ のある有限拡大体 $\mathsf{Z}$ に属する。$\mathfrak K_{\mathsf{Z}}$ の中で $\mathfrak a'=z_1\mathsf{Z}+\cdots+z_r\mathsf{Z}$ は両側イデアルである。$\mathfrak K_{\mathsf{Z}}$ は両側単純なので、$\mathfrak a'$ は零または $\mathfrak K_{\mathsf{Z}}$ に等しい。したがって $\mathfrak a$ は零または $\mathfrak K_\Omega$ に等しく、$\mathfrak K_\Omega$ は両側単純である。

この定理により、$\mathfrak K_\Omega$ はその中心 $\Omega$ のある拡大体上の行列環である。この体は $\mathfrak K_\Omega$ の部分環として $\Omega$ 上有限次数でなければならず、したがって $\Omega$ 自身に等しい。

\textbf{定理 3.} \emph{$\mathfrak K_\Omega$ は $\Omega$ 上の行列環である：}
\[
\mathfrak K_\Omega=\sum\Omega c_{ik}.
\]
\emph{したがって、中心 $\mathsf P$ 上の $\mathfrak K$ の次数は平方数 $m=t^2$ である。}

$\mathfrak K_\Omega$ が分解する単純右（または左）イデアルの個数 $t$ を、$\mathfrak K$ の絶対成分数と呼ぶ。

$\mathfrak K_\Omega$ が $t$ 個の単純イデアルへ分解するという $\Omega$ の性質は、$\mathsf P$ のある有限拡大 $\mathsf{Z}$ によってすでに得られる。たとえば $c_{ik}$ の係数を $\mathsf P$ に添加して得られる $\mathsf{Z}$ がそうである。次にこれらの分解体を調べる。

\textbf{定義。} $\mathsf P$ の有限代数拡大体 $Z$ を $\mathfrak K$ の\emph{分解体}というのは、$\mathfrak K_{\mathsf{Z}}$ が分解する単純右イデアルの個数がすでに絶対成分数 $t$ に等しいときである。\srcfn{1}{9 頁の分解体の概念と比較せよ。}

\textbf{定理 4.} \emph{$\mathsf{Z}$ が $\mathfrak K$ の分解体ならば、$\mathsf{Z}$ は $\mathfrak K_{\mathsf{Z}}$ の右イデアルの自己同型体である。逆に、$\mathfrak K_{\mathsf{Z}}$ の右イデアルの自己同型体が $\mathsf{Z}$ ならば、$\mathsf{Z}$ は分解体である。}

\emph{証明。} 1. $\mathsf{Z}$ を分解体、$\mathfrak T$ を $\mathfrak K_{\mathsf{Z}}$ の右イデアルの自己同型体とする。すなわち
\[
\mathfrak K_{\mathsf{Z}}=\sum\mathfrak T c_{ik}.
\]
ここから
\[
\mathfrak K_\Omega=\sum\mathfrak T_\Omega c_{ik}
\]
が従う。したがって $\Omega$ 上の $\mathfrak K_\Omega$ の階数は $t^2$ と $\mathfrak T_\Omega$ の階数との積である。他方、その階数は $t^2$ に等しいから、$\mathfrak T_\Omega$、したがって $\mathfrak T$ の階数は $1$ であり、$\mathfrak T=\mathsf{Z}$ である。

\textbf{2.} $\mathsf{Z}$ が $\mathfrak K_{\mathsf{Z}}$ の単純イデアルの自己同型体ならば
\[
\mathfrak K_{\mathsf{Z}}=\sum\mathsf Zc_{ik},
\]
したがって
\[
\mathfrak K_\Omega=\sum\Omega c_{ik}.
\]
ゆえに $c_{ik}$ の個数は $t^2$ でなければならない。証明された。

\textbf{定理 5.} \emph{$\mathfrak K$ の分解体 $Z$ の $\mathfrak K$ における既約表現 $\mathfrak Z$ は、行列環 $\mathfrak K_r$ の極大可換部分体である（22 頁の定理 2 により既約表現類は一つだけであり、$r$ は $\mathsf{Z}$ の既約表現の次数である）。}

\emph{逆に、$\mathsf P$ の可換拡大体 $Z$ の既約表現 $\mathfrak Z$ が $\mathfrak K_r$ の極大可換部分体ならば、$Z$ は分解体である。}

\emph{証明。} 1. $\mathsf P$ 上の分解体 $\mathsf{Z}$ の次数を $n$ とすると、$\mathsf{Z}$ の既約表現 $\mathfrak Z$ の次数 $r$ は $n/t$ である（22 頁）。$\mathfrak Z\subseteqq\mathfrak Z^*\subseteqq\mathfrak K_r$ とし、$\mathfrak Z^*$ を $\mathsf P$ 上次数 $n^*\geqq n$ の極大可換体とする。$\mathfrak Z^*=\mathfrak Z$ を示せばよいが、その代わりに $n^*=n$ を示せば十分である。$Z$ の拡大体で $\mathfrak Z^*$ と同型な $\mathsf{Z}^*$ の既約表現の次数を $r^*$ とすると、再び $n^*=r^*t$ である。他方、$\mathfrak Z^*$ の中には $\mathsf{Z}^*$ の $r$ 次表現があるので、$r$ は $r^*$ で割り切れ、$r^*\leqq r$ である。したがって
\[
n\leqq n^*=r^*t\leqq rt=n \qquad\text{より}\qquad n^*=n.
\]

2. $\mathsf P$ の可換拡大体 $\mathsf{Z}$ の $\mathfrak K$ における $r$ 次既約表現 $\mathfrak Z$ が、$\mathfrak K_r$ の極大可換部分体であるとする。$\mathfrak K_{\mathsf{Z}}$ の単純イデアルの自己同型体を $\mathfrak T$ とし、
\[
\mathfrak K_{\mathsf{Z}}=\sum_{1}^{t'}\mathfrak T c_{ik}
\]
とする。$\mathfrak T$ は $\mathsf{Z}$ を含む。$\mathfrak T=\mathsf{Z}$ を示さなければならない。$\mathfrak K_r$ の単純右イデアル $\mathfrak r$ は、$\mathsf{Z}$ の $\mathfrak K$ における表現加群であるだけでなく、$\mathfrak T$ の表現加群でもある。実際、
\[
\mathfrak r=\mathfrak T c_{i1}+\cdots+\mathfrak T c_{ik}
\]
であるから、$\mathfrak T\mathfrak r=\mathfrak r$ であり、$\mathfrak r$ は $\mathfrak T$-左加群かつ $\mathfrak K$-右加群である。

したがって $\mathfrak K_r$ には $\mathfrak T$ と同型で $\mathfrak Z$ を含む部分体 $\overline{\mathfrak T}$ がある。$\mathfrak T$ が $\mathsf{Z}$ と異なり、したがって $\overline{\mathfrak T}$ が $\mathfrak Z$ と異なるならば、$\mathfrak Z$ に属さない $\overline{\mathfrak T}$ の元 $\alpha$ を $\mathfrak Z$ に添加すると、$\mathfrak K_r$ 内で $\mathfrak Z$ の真の可換拡大体が得られる。これは $\mathfrak Z$ の極大性に反する。証明された。

定理 5 から直ちに次が従う。

\textbf{定理 6.} \emph{$\mathfrak K$ の極大可換部分体は $\mathsf P$ 上次数 $t$ であり、分解体である。\srcfn{1}{各分解体が $\mathfrak K$ の極大可換部分体を含むという推測は誤りである。その性質をもたない分解体もありうる。そのとき対応する $r$ は必ず $1$ より大きい。R. Brauer and E. Noether, ``Über minimale Zerfällungskörper irreduzibler Darstellungen'', Sitzungsberichte d. preuss. Akademie der Wissenschaften 1927 XXXII, 221 を参照。}}

次に $\mathfrak K_r$ の部分環の同型に関する定理を導く。非可換体 $\mathfrak K$ のガロア群の構造と、それに関連する諸事実の精密な理解はこの定理に基づく。

\textbf{定理 7.} \emph{$\mathfrak S_1,\mathfrak S_2$ を、$\mathsf P$ を含む、$\mathfrak K$ 上の行列環 $\mathfrak K_r$ の両側単純部分環とし、両者の間に $\mathsf P$ を元ごとに固定する同型 $\mathfrak S_1\cong\mathfrak S_2$ があるとする。このとき、その同型は $\mathfrak K_r$ の正則元 $\tau$ による変換である。}

\emph{$\sigma_1\in\mathfrak S_1$ と $\sigma_2\in\mathfrak S_2$ が対応するならば、}
\[
\tau^{-1}\sigma_1\tau=\sigma_2.
\]

\emph{証明。} $\mathfrak S_1$ と $\mathfrak S_2$ の双方に相反同型な環 $\Sigma$ を構成し、その中の $\mathsf P$ と同型な部分体を $\mathsf P$ と同一視する。すると $\mathfrak S_1,\mathfrak S_2$ は $\Sigma$ の $\mathfrak K$ における同じ次数 $r$ の二つの相反表現である。22 頁の定理 2 により既約相反表現類は一つだけであり、したがって $\Sigma$ の $\mathfrak K$ における所定の次数の表現類も一つだけである（既約類の倍数）。ゆえに $\mathfrak S_1$ と $\mathfrak S_2$ は同じ類に属し、変換によって互いに移り合う。証明された。

定理 7 から直ちに次が従う。

\textbf{定理 8.} \emph{$\mathsf P$ を元ごとに固定する $\mathfrak K$ の自己同型群、すなわち中心に関する $\mathfrak K$ のガロア群 $\mathfrak G$ は、$\mathfrak K$ の内部自己同型群である。}

定理 7 で $r=1$、$\mathfrak S_1=\mathfrak S_2=\mathfrak K$ と置けばよい。

次の二定理は、ここまでに展開した一般理論が個別の問題に有用であることを示す。

\textbf{定理 9.} \emph{中心が実数体 $\mathsf{R}$ である有限次数の非可換体は、四元数体だけである。}

\emph{証明。} $\mathfrak K$ を中心 $\mathsf{R}$ の非可換体（$\mathsf{R}$ 上有限階数）とする。定理 6 により、$R$ 上の $\mathfrak K$ の次数は、$\mathfrak K$ の極大可換部分体の次数 $t$ の平方である。この部分体は複素数体と同型な体しかありえないから、$t=2$ であり、$\mathfrak K$ は $\mathsf{R}$ 上次数 $4$ である。$\mathfrak Z_1$ を $\mathfrak K$ の極大可換部分体とし、$i^2=-1$ として $\mathfrak Z_1=\mathsf{R}(i)$ と置ける。$\mathfrak Z_2=\mathsf{R}(-i)$ も $\mathfrak Z_1$ と同型であり、集合としては $\mathfrak Z_1$ と一致する体で、その同型は $i\mapsto-i$ で与えられる。したがって定理 7 により、
\[
j^{*-1}ij^*=-i
\]
を満たす $j^*\in\mathfrak K$ が存在する。

$j^*$ は $\mathfrak Z_1$ に属しえないので、$\mathfrak Z_1$ 上の $\mathfrak Z_1(j^*)$ の階数は少なくとも $2$ である。したがって $\mathfrak Z_1(j^*)$ は $\mathsf{R}$ 上階数 $4$ で、$\mathfrak K$ に等しく、
\[
\mathfrak K=\mathsf R+\mathsf Ri+\mathsf Rj^*+\mathsf Rij^*
\]
と書ける。

\[
\begin{aligned}
j^{*-2}1j^{*2}&=1\\
j^{*-2}ij^{*2}&=i\\
j^{*-2}j^*j^{*2}&=j^*\\
j^{*-2}ij^*j^{*2}&=ij^*
\end{aligned}
\]
であるから、$j^{*2}$ は $\mathfrak K$ のすべての元と可換で、中心 $\mathsf{R}$ に属し、実数 $r$ である。$r$ は正ではありえない。正ならば、可換体 $\mathsf{R}(j^*)$ の多項式 $x^2-r$ が $+\sqrt r,-\sqrt r,+j^*,-j^*$ という四つの根をもつことになるからである。したがって $j^{*2}=-m^2$、$m\in \mathsf{R}$ である。$j=j^*m^{-1}$ と置けば $j^2=-1$ となり、$1,i,j,k=ij$ は四元数単位である。証明された。

この証明は、$\mathsf{R}$ の代わりに任意の代数的実閉体について明らかに成り立つ。

\textbf{定理 10（Wedderburn の定理）。} \emph{有限個の元からなる体はすべて可換である。}

\emph{証明。} $\mathfrak K$ を有限体、$\mathsf P$ をその中心とし、$\mathsf P$ 上の $\mathfrak K$ の次数を $t^2$ とする。$\mathfrak K$ の極大可換部分体はすべて $\mathsf P$ 上次数 $t$ なので、同じ個数の元をもつ。したがってガロア体の理論の周知の定理（Moore の定理）により互いに同型である。さらに各ガロア体はその素体のガロア拡大なので、同型は $\mathsf P$ の各元を固定するように選べる。定理 7 により、$\mathfrak K$ の極大可換部分体はすべて、一つの $\mathfrak Z$ から変換によって得られる。一方、$\mathfrak K$ はその極大可換部分体すべての合併である（$\mathsf P$ に\emph{一つ}の元を添加すれば可換体が生じる）。ここで $\mathfrak K$ の乗法群 $\mathfrak K^*$（零を除いた $\mathfrak K$）だけを考えると、$\mathfrak K^*$ は部分群 $\mathfrak Z^*$ と共役な部分群すべての合併である。しかし $t>1$ ならば矛盾が生じる。$\mathfrak K^*$ が $M$ 個、$\mathfrak Z^*$ が $N$ 個の元をもち、$L$ を $\mathfrak K^*$ における $\mathfrak Z^*$ の指数とすると、$M=NL$ であり、相異なる共役部分群の個数は高々 $L$ である。この個数が $L$ に等しいとしても、それらの合併が群全体を尽くすには、どの二つも共通の元をもってはならない。しかしそれらはすべて単位元を共有する。

\subsection*{§ 21. 非可換体のガロア理論}\label{die-galoissche-theorie-der-nichtkommutativen-koerper}

$\mathsf{G}$ を中心 $\mathsf P$ に関する $\mathfrak K$ のガロア群とする。定理 8 により $G$ は $\overline{\mathsf G}=\mathfrak K^*/\mathsf P^*$ と同型であり、この同型を固定する。$G$ の部分群 $H$ を\emph{閉}と呼ぶのは、同型 $\mathsf{G}\cong\mathfrak K^*/\mathsf P^*$ によって対応する $\mathfrak K^*$ の部分群 $\mathfrak H^*$ に\emph{零}を加えたものが体となるときである。

\textbf{定理 11　ガロア理論の基本定理。} \emph{$G$ の閉部分群 $H$ と $\mathsf P$ と $\mathfrak K$ の中間体 $\mathfrak S$ との間には一対一対応がある。$\mathsf{H}$ と $\mathfrak S$ が対応するとき、$\mathsf{H}$ は $\mathfrak S$ の不変群であり、$\mathfrak S$ は $\mathsf{H}$ の不変体である。}

\emph{証明。} 定理を三つの部分命題に分ける。
\begin{enumerate}
\item 体 $\mathfrak S$ の不変群 $H$ は閉である。
\item $\mathfrak S$ の不変群が $\mathsf{H}$ ならば、$\mathfrak S$ は $\mathsf{H}$ の不変体である。
\item $\mathsf{H}$ の不変体が $\mathfrak S$ ならば、$\mathsf{H}$ は $\mathfrak S$ の不変群である。
\end{enumerate}

\begin{enumerate}
\item 同型 $\mathsf{G}\cong \mathfrak K^*/\mathsf P^*$ により、体 $\mathfrak S$ の不変群 $\mathsf{H}$ には、$\mathfrak S$ のすべての元と可換な零でない $\mathfrak K$ の元全体 $\mathfrak H^*$ が対応する。これらに $0$ を加えれば明らかに体となる。したがって $\mathsf{H}$ は閉である。

\item 第三の部分命題を第二へ帰着させる。体 $\mathfrak S$ には $\mathsf{G}$ の部分群 $\mathsf{S}$、すなわち $\mathfrak S^*/\mathsf P^*$ に対応する部分群が属する。同様に群 $\mathsf{H}$ には $\mathfrak K$ の部分体 $\mathfrak H$ が属し、$\mathsf{H}\cong \mathfrak H^*/\mathsf P^*$ である。$\mathfrak S$ が $\mathsf{H}$ の不変体ならば、$\mathfrak S$ は $\mathfrak H$ の各元と可換な $\mathfrak K$ の元すべてからなる。したがって、$\mathsf{S}$ は $\mathfrak H$ の不変群であるともいえる。同様に、$\mathsf{H}$ が $\mathfrak S$ の不変群ならば、$\mathfrak H$ は $\mathsf{S}$ の不変体である。ゆえに第三の部分命題は次のようにも述べられる。

\item[$3'$.] $\mathfrak H$ の不変群が $\mathsf{S}$ ならば、$\mathfrak H$ は $\mathsf{S}$ の不変体である。これはまさに、なお証明すべき第二の部分命題である。

\item[3.] この証明は若干の小さな相違を除いて可換の場合と同様に進む。
\end{enumerate}

$\mathfrak K$ と相反同型な体 $\mathsf{K}$ を作り、その中心を $\mathsf P$ と同一視する。$\mathfrak K$ の自己同型の代わりに、$\mathfrak S$ から $\mathsf{K}$ への相反同型を考えられる。

可換の場合に用いた延長定理に類似する補題を再び証明する。$\mathfrak S,\mathfrak T$ を $\mathsf P$ と $\mathfrak K$ の中間体とし、
\[
\mathsf P \subseteqq \mathfrak S \subseteqq \mathfrak T \subseteqq \mathfrak K
\]
とする。$\mathfrak S$ から $\mathsf{K}$ の部分体への各相反同型（$\mathfrak S$ の $\mathsf{K}$ における一次相反表現）は、$\mathfrak S$ 上の $\mathfrak T$ の次数 $s$ と少なくとも同じ数の仕方で、$\mathfrak T$ の相反同型へ延長できると主張する。

体 $\mathsf P,\mathfrak S,\mathfrak T,\mathfrak K$ を $\mathsf{K}$ によって拡大する：
\[
\mathsf{K}\subseteqq \mathfrak S_{\mathsf{K}}\subseteqq \mathfrak T_{\mathsf{K}}\subseteqq \mathfrak K_{\mathsf{K}}.
\]
20 頁の定理 1 により、各 $\mathfrak S_{\mathsf{K}}$ は両側単純かつ完全可約であり、互いに作用素同型な単純左イデアルは、$\mathfrak S$ の $\mathsf{K}$ におけるただ一つの既約表現類の表現加群である。この表現類は一次である。$\mathsf{K}$ の構成により、$\mathsf{K}$ の中には $\mathfrak S$ と相反同型な体が存在するからである。したがって単純左イデアルはすべて階数一であり、$\mathfrak S_{\mathsf{K}}$ はその $\mathsf P$ 上の次数と同じ個数の左イデアルの和となる：
\[
\mathfrak S_{\mathsf{K}}=\mathsf K E_1+\cdots+\mathsf K E_p,\qquad \mathsf K E_\nu=\mathfrak S_{\mathsf{K}} E_\nu.
\]
すると
\[
\begin{aligned}
\mathfrak T_{\mathsf{K}}&=\mathfrak T\cdot \mathsf{K}=\mathfrak T\cdot \mathfrak S_{\mathsf{K}},\\
\mathfrak T_{\mathsf{K}}&=\mathfrak T\mathfrak S_{\mathsf{K}} E_1+\cdots+\mathfrak T\mathfrak S_{\mathsf{K}} E_p
              =\mathfrak T_{\mathsf{K}} E_1+\cdots+\mathfrak T_{\mathsf{K}} E_p.
\end{aligned}
\]
である。各イデアル $\mathfrak T_{\mathsf{K}}\cdot E_\nu$ は $s$ 個の単純イデアルへ分解する。実際、
\[
[\mathfrak T_{\mathsf{K}} E_\nu:\mathsf{K}]=[\mathfrak T_{\mathsf{K}} E_\nu:\mathsf K E_\nu]=[\mathfrak T_{\mathsf{K}} E_\nu:\mathfrak S_{\mathsf{K}} E_\nu]\leqq[\mathfrak T_{\mathsf{K}}:\mathfrak S_{\mathsf{K}}]=[\mathfrak T:\mathfrak S]=s
\]
であり、
\[
\sum [\mathfrak T_{\mathsf{K}} E_\nu:\mathsf{K}]=sp
\]
だから、
\[
[\mathfrak T_{\mathsf{K}} E_\nu:\mathsf{K}]=s,
\qquad
\mathfrak T_{\mathsf{K}} E_\nu=\mathsf K e_1^{(\nu)}+\cdots+\mathsf K e_s^{(\nu)}.
\]

$\mathsf K e_i^{(\nu)}$ が与える $\mathfrak T$ の表現は、$\mathsf K E_\nu$ が与える $\mathfrak S$ の表現の延長である。実際、$\alpha\in\mathfrak S$ ならば
\[
\alpha E_\nu=\sigma E_\nu,\qquad
\sum \alpha e_i^{(\nu)}=\sigma\sum e_i^{(\nu)},\qquad
\alpha e_i^{(\nu)}=\sigma e_i^{(\nu)}.
\]

残るのは、$\mathsf K e_i^{(\nu)}$ が与える表現がすべて相異なることの証明だけである。可換の場合には、それらが異なる表現類に属したので明らかであった。ここでは逆にすべて同じ表現類に属するため、表現そのものを調べなければならない。$\mathsf K e_i^{(\nu)}$ が与える $\mathfrak T_{\mathsf{K}}$ 全体の $\mathsf{K}$-準同型写像を考える。それらは $\mathfrak T$ の表現がわかれば完全に定まる。写像全体を構成するには、基底元がどのように表現されるかを知れば十分だからである。したがって、$\mathsf K e_i^{(\nu)}$ が与える $\mathfrak T_{\mathsf{K}}$ の写像がすべて異なることを示せばよい。$e_i^{(\nu)}$ を一度は表現加群 $\mathsf K e_i^{(\nu)}$ により、別には $\mathsf K e_j^{(\nu)}$ により表現すると、
\[
e_i^{(\nu)2}=e_i^{(\nu)}
\]
により前者では $1$ が、
\[
e_i^{(\nu)}e_j^{(\nu)}=0
\]
により後者では $0$ が、$e_i^{(\nu)}$ を表す元となる。したがって二つの表現は実際に異なる（全く同じ議論は可換の場合にも使えるが、そこでは不要である）。

ここからは 10 頁および 16 頁と全く同じように論じる。$\mathsf{H}$ を $\mathfrak S$ の不変群、$\mathfrak T$ を $\mathsf{H}$ の不変体とする。$\mathsf{H}$ の元は、$\mathfrak S$ の恒等自己同型を $\mathfrak K$ 全体へ延長したものすべてである。直前の補題により、その中には $\mathfrak T$ 上で相異なるものが少なくとも $s$ 個あるから、$s=1$、$\mathfrak T=\mathfrak S$ でなければならない。証明された。

